\subsection{System Overview}
\label{sec:system-overview}

Figure~\ref{fig:system-overview} separates applications, the control
plane, and inference engines. Its numbered badges identify mechanisms
discussed below; they do not specify execution order.

\begin{figure*}[t]
  \centering
  \includegraphics[width=\textwidth]{figures/system-overview-v11.drawio.pdf}
  \caption{AegisKV overview. The four annotations identify Workflow Tracker
    (1), KV Placement Planner (2), Recovery Manager (3), and Adapter (4).
    DP0 and DP1 have equivalent components; badge 4 is shown once.
    The Planner selects tiers R/D/P (recompute, DRAM, or Peer competition)
    and layer/block sets. Outlined local blocks illustrate a Peer selection
    in early layers; the gold final column in Local KV is the unsealed tail.
    Black arrows dispatch requests; gray dashed arrows carry control
    information; red dashed arrows submit retries. Thin gray lines denote
    DAG dependencies or local read/restore interfaces. Blue paths transfer Peer KV;
    gold paths independently copy to and restore from DRAM. Control arrows
    show representative instances. The retry fork represents concurrent
    submissions. On takeover, the target reads its Peer KV locally.
    Block counts, bar lengths, and memory proportions are schematic.}
  \label{fig:system-overview}
  \Description{Three horizontal tiers show applications, control, and inference.
    A workflow DAG feeds a Tracker with per-request slack bars. The Planner
    selects protection tiers and layer/block sets; the Recovery Manager forks
    concurrent retries. Equivalent DP0 and DP1 engines each contain an Adapter,
    a layer/block grid in Local KV, a smaller Peer KV region holding the other
    DP's replicas, and a Connector. Adapters invoke Connectors; Peer transfers
    and independent DRAM copy/restore paths bypass the control plane.}
\end{figure*}

\paragraph{Applications.}
The agent runtime submits LLM requests with workflow and node IDs and
declares dependency edges. Static DAGs can be supplied in advance;
new nodes are reported as they are spawned. Applications need
not estimate slack.

\paragraph{Control plane.}
The Router dispatches normal requests using load and KV locality.
The AegisKV Controller owns workflow state, protection plans, and recovery
coordination. Workflow Tracker (1) produces slack snapshots
(Section~\ref{sec:workflow-protection}). KV Placement Planner (2) selects
tiers, destinations, layer/block sets, and transfer plans from these
snapshots, calibrated profiles, recent block rankings, and available
capacity (Sections~\ref{sec:kv-placement} and~\ref{sec:capacity-allocation}).
Recovery Manager (3) combines membership, engine/GPU progress, and error
reports to initiate takeover (Section~\ref{sec:parallel-recovery}).
The KV directory and commit metadata are shared controller state that
must survive the source engine. Calibration tables are configuration,
not a separate online service.

\paragraph{Inference engines.}
DP0 and DP1 each have an Adapter (4), a Connector, and equivalent memory
regions. Local KV serves that DP's requests; a smaller, reclaimable Peer
region holds another DP's backups. The Adapter reports selector rankings
and capacity, manages sealing and commits, and implements reclamation and
recovery loading (Section~\ref{sec:kv-runtime}). The Connector executes
transfers selected by the Planner. KV moves directly between engines and
storage, without passing through the Controller.

By default the Planner refreshes its plan every 1\,s; request arrival,
completion, or a significant slack change can trigger an earlier refresh.
The Adapter reconciles the desired objects with committed copies and
issues only incremental transfers. Peer and DRAM paths commit independently.
Plans, rankings, and commit state remain available after source failure.

\paragraph{Predetermined takeover.}
A Peer destination is chosen before failure so that the target can read
the replica locally. A simple mapping is $p(g)=(g+1)\bmod n$ for $n\ge2$
DPs. Each DP then backs up one source DP; tensor-parallel deployments
treat a whole TP group as a DP and restore compatible shards.
Spreading a source's requests over several recovery targets requires
binding requests and replicas before failure. An arbitrary post-failure
target change would lose the no-second-copy property.

\begin{table}[t]
  \caption{Decision timescales. Block identities change more frequently
    than calibrated hot-set sizes.}
  \label{tab:design-timescales}
  \centering
  \begin{tabularx}{\columnwidth}{@{}lX@{}}
    \toprule
    Timescale & Decisions and ownership \\
    \midrule
    Calibration & Block size, $K_l$, costs, and effective bandwidth;
      stored as Planner configuration. \\
    Placement window & Tiers, slack order, prefix $L$, footprint $F$,
      and budget $B$; Tracker and Planner. \\
    Ranking window & Protected block identities; Adapter counters
      feed Planner selection. \\
    Execution & Seal, transfer, commit, reclaim, and demand-fetch;
      Adapter and Connector. \\
    \bottomrule
  \end{tabularx}
\end{table}
